\section*{Введение}
\addcontentsline{toc}{section}{Введение}

Объектом исследования, как следует из названия, являются цифровые двойники, а конкретно --- представление объектов реального мира в форме, которая позволит добиться их фотореалистичной визуализации.

Визуальное представление мира строится на свете, поэтому в первую очередь исследованию стоит подвергнуть симуляцию поведения света, которое (поведение) разделяется два случая:
\begin{itemize}
	\item На границе раздела двух сред с разным коэффициентом приломления часть света отражается, часть приломляется; математической моделью этого процесса в рендеринге\footnote{Под словом рендеринг здесь и далее понимается симуляция света (различной степени физической корректности) для получения картинки на сенсоре, которым в компьютерах выступает монитор} выступают двунаправленные функции распределения рассеивания (BSDF), за отражение и приломление отвечают соотвественно двунаправленные функции отражения и пропускания (BRDF, BTDF);
	\item В участвующих средах свет непосредственно взаимодействует и подвергается поглощению и рассеиванию, это описывается уравнением переноса излучения (RTE).
\end{itemize}

В таком случае цифровым двойником можно называть описывающие объект оптические свойства, которые понимает реализующее симуляцию поведения света программное обеспечение --- движок рендеринга.

В настоящей работе предлагается ограничить область рассмотрения до медицинских потребностей, где ведутся активные исследования как в части изучения оптических свойств \cite{Jacques_2013}, так и в разработке движков рендеринга \cite{kroes2012exposure}.

Из последнего удобно подчерпнуть идеи интерпретации медицинских данных, которые получаются в основном из КТ (выдаёт единицы Хаундсфилда, HU) и МРТ в виде скалярной воксельной сетки. Поэтому перед тем, как загонять их в движок, требуется некоторая подготовка. Для этого вводится функция переноса, чаще всего $\mathrm{HU} \to (RGBA)$, а градиент $\nabla\mathrm{HU}$ используется для выделения границ органов.

Иногда удается сопоставлять и более сложные оптические свойства. Хорошим примером является работа \cite{verleker2014optical}, в которой из КТ и априорных знаниях о том, каким значениям HU соответствуют те или иные ткани, извлекают коэффициенты поглощения, рассеивания и анизотропию (описывающую распределение рассеивания света), по которым симулируется поведение света, чтобы узнать попадает ли он в нужные зоны головного мозга.

Имеет место предположение, что такое дотошное следование законам физики излишне, и возможно упрощение цифровых двойников до модели 3DGS, которое позволит достигать всё тех же целей (фотореалистиная физуалиция) но с улучшением скорости рендеринга.